% ============================================================
%  ECON306: W11D2 — Human Capital vs. Signalling: Evidence;
%          Sheepskin Effects
%  Date: Thursday 1 October 2026
%  SOURCE: old w13_week13 (evidence section) from FULL deck
%    (lines 3927–3976) + Course Map W11D2 LOs.
% ============================================================

\documentclass[aspectratio=169, 12pt]{beamer}

% ============================================================
%  ECON306: Shared Preamble — The Economics of Health and Education
%  University of Otago — Semester 2, 2026
%
%  SHARED BY ALL ECON306 DAY DECKS.
%  Input this file immediately after \documentclass in each deck:
%    % ============================================================
%  ECON306: Shared Preamble — The Economics of Health and Education
%  University of Otago — Semester 2, 2026
%
%  SHARED BY ALL ECON306 DAY DECKS.
%  Input this file immediately after \documentclass in each deck:
%    % ============================================================
%  ECON306: Shared Preamble — The Economics of Health and Education
%  University of Otago — Semester 2, 2026
%
%  SHARED BY ALL ECON306 DAY DECKS.
%  Input this file immediately after \documentclass in each deck:
%    \input{../Preambles/econ306_preamble.tex}
%  Do NOT include \documentclass or per-deck metadata here.
% ============================================================

% ─── Theme & Colours ────────────────────────────────────────
\usetheme{default}
\usecolortheme{default}
\usefonttheme{professionalfonts}
\definecolor{OtagoBlue}{HTML}{003057}
\definecolor{OtagoGold}{HTML}{A8923A}
\definecolor{TextGrey}{HTML}{3A3A3A}
\definecolor{AccentBlue}{HTML}{2B6CB0}
\definecolor{LightBg}{HTML}{F7F9FC}

\setbeamercolor{frametitle}{fg=OtagoBlue, bg=white}
\setbeamercolor{title}{fg=OtagoBlue}
\setbeamercolor{subtitle}{fg=OtagoGold}
\setbeamercolor{author}{fg=TextGrey}
\setbeamercolor{date}{fg=TextGrey}
\setbeamercolor{institute}{fg=TextGrey}
\setbeamercolor{normal text}{fg=TextGrey, bg=white}
\setbeamercolor{item}{fg=OtagoBlue}
\setbeamercolor{subitem}{fg=AccentBlue}
\setbeamercolor{block title}{bg=OtagoBlue!12, fg=OtagoBlue}
\setbeamercolor{block body}{bg=OtagoBlue!5, fg=TextGrey}
\setbeamercolor{alertblock title}{bg=OtagoGold!20, fg=OtagoBlue}
\setbeamercolor{alertblock body}{bg=OtagoGold!8, fg=TextGrey}
\setbeamercolor{alerted text}{fg=OtagoGold!80!black}
\setbeamercolor{structure}{fg=OtagoBlue}

\setbeamertemplate{frametitle}{%
	\vspace{0.35em}%
	\textbf{\insertframetitle}\par%
	\vspace{0.05em}%
	\textcolor{OtagoGold}{\rule{\textwidth}{0.6pt}}}

\setbeamertemplate{navigation symbols}{}
\setbeamertemplate{footline}{%
	\leavevmode\hbox{%
		\begin{beamercolorbox}[wd=.40\paperwidth,ht=2.5ex,dp=1.5ex,leftskip=.3cm]{author in head/foot}%
			\textcolor{TextGrey}{\footnotesize ECON306 \textbullet\ University of Otago}%
		\end{beamercolorbox}%
		\begin{beamercolorbox}[wd=.47\paperwidth,ht=2.5ex,dp=1.5ex,center]{title in head/foot}%
			\textcolor{TextGrey}{\footnotesize \insertshorttitle}%
		\end{beamercolorbox}%
		\begin{beamercolorbox}[wd=.13\paperwidth,ht=2.5ex,dp=1.5ex,rightskip=.3cm,right]{date in head/foot}%
			\textcolor{TextGrey}{\footnotesize \insertframenumber\,/\,\inserttotalframenumber}%
	\end{beamercolorbox}}\vskip0pt}

\setbeamertemplate{itemize item}{\textcolor{OtagoGold}{\small$\blacktriangleright$}}
\setbeamertemplate{itemize subitem}{\textcolor{AccentBlue}{\small$\triangleright$}}
\setbeamertemplate{enumerate item}{\textcolor{OtagoBlue}{\small\bfseries\insertenumlabel.}}

\setbeamertemplate{title page}{%
	\begin{tikzpicture}[remember picture,overlay]
		\fill[OtagoBlue] (current page.north west)
		rectangle ([yshift=-0.38\paperheight]current page.north east);
	\end{tikzpicture}
	\vspace{0.06\paperheight}
	\begin{beamercolorbox}[wd=\textwidth, sep=0pt]{title}
		{\color{white}\Large\textbf{\inserttitle}}\\[0.4em]
		{\color{OtagoGold}\large\insertsubtitle}
	\end{beamercolorbox}
	\vspace{1.4em}
	\begin{beamercolorbox}[wd=\textwidth, sep=0pt]{author}
		{\color{TextGrey}\normalsize\insertauthor}\\[0.2em]
		{\color{TextGrey}\small\insertinstitute}\\[0.2em]
		{\color{TextGrey}\small\insertdate}
\end{beamercolorbox}}

% ─── Packages ──────────────────────────────────────────────
\usepackage{tikz}
\usepackage{booktabs}
\usepackage{array}
\usepackage{graphicx}
\usepackage{hyperref}
\usepackage{amsmath}
\usepackage{amssymb}
\usepackage{colortbl}
\usepackage{multirow}
\usepackage{xcolor}
\hypersetup{colorlinks=true, linkcolor=AccentBlue, urlcolor=AccentBlue,
	citecolor=AccentBlue, pdfauthor={Austin Henderson},
	pdftitle={ECON306 S2 2026}}

% ─── Images folder ─────────────────────────────────────────
%  Place all slide images in the 'images/' subfolder next to each deck.
%  Usage: \includegraphics[width=0.85\textwidth]{figure_name}
\graphicspath{{images/}}

% ─── Reusable macros ───────────────────────────────────────
% Recap slide at start of each lecture
\newcommand{\recapslide}[1]{%
	\begin{frame}{Recap: What Did We Cover Last Time?}
		\begin{block}{Key points from last lecture}
			\begin{itemize}#1\end{itemize}
		\end{block}
\end{frame}}

% Plan slide at start of each lecture
\newcommand{\planslide}[2]{%
	\begin{frame}{Today's Plan — #1}
		\begin{columns}
			\column{0.6\textwidth}
			\begin{itemize}#2\end{itemize}
			\column{0.38\textwidth}
			\begin{alertblock}{Reading}
				\footnotesize See Brightspace (Aoroa) for required readings before this lecture.
			\end{alertblock}
		\end{columns}
\end{frame}}

% Summary slide at end of each lecture
\newcommand{\summaryside}[1]{%
	\begin{frame}{Lecture Summary}
		\begin{exampleblock}{Key takeaways}
			\begin{itemize}#1\end{itemize}
		\end{exampleblock}
		\vspace{0.5em}
		{\footnotesize\textcolor{OtagoGold}{\textbf{Brightspace (Aoroa):}} slides, readings, and any announcements posted there.}
\end{frame}}

% NZ spotlight box
\newcommand{\nzbox}[2]{%
	\begin{alertblock}{\includegraphics[height=0.8em]{nz_flag_icon}\hspace{0.3em} NZ Spotlight: #1}
		#2
\end{alertblock}}

% Tutorial callout
\newcommand{\tutorialbox}[1]{%
	\begin{block}{\textbf{Tutorial This Week}}
		#1
\end{block}}

% ─── Section dividers ──────────────────────────────────────
\AtBeginSection[]{
	\begin{frame}[plain]
		\begin{tikzpicture}[remember picture,overlay]
			\fill[OtagoBlue!95] (current page.north west)
			rectangle (current page.south east);
			\node[anchor=center, text=white, font=\Large\bfseries,
			text width=0.9\paperwidth, align=center]
			at (current page.center) {\insertsectionhead};
			\node[anchor=south, text=OtagoGold, font=\normalsize,
			text width=0.9\paperwidth, align=center]
			at ([yshift=2.5cm]current page.south) {ECON306 — University of Otago};
		\end{tikzpicture}
\end{frame}}

%  Do NOT include \documentclass or per-deck metadata here.
% ============================================================

% ─── Theme & Colours ────────────────────────────────────────
\usetheme{default}
\usecolortheme{default}
\usefonttheme{professionalfonts}
\definecolor{OtagoBlue}{HTML}{003057}
\definecolor{OtagoGold}{HTML}{A8923A}
\definecolor{TextGrey}{HTML}{3A3A3A}
\definecolor{AccentBlue}{HTML}{2B6CB0}
\definecolor{LightBg}{HTML}{F7F9FC}

\setbeamercolor{frametitle}{fg=OtagoBlue, bg=white}
\setbeamercolor{title}{fg=OtagoBlue}
\setbeamercolor{subtitle}{fg=OtagoGold}
\setbeamercolor{author}{fg=TextGrey}
\setbeamercolor{date}{fg=TextGrey}
\setbeamercolor{institute}{fg=TextGrey}
\setbeamercolor{normal text}{fg=TextGrey, bg=white}
\setbeamercolor{item}{fg=OtagoBlue}
\setbeamercolor{subitem}{fg=AccentBlue}
\setbeamercolor{block title}{bg=OtagoBlue!12, fg=OtagoBlue}
\setbeamercolor{block body}{bg=OtagoBlue!5, fg=TextGrey}
\setbeamercolor{alertblock title}{bg=OtagoGold!20, fg=OtagoBlue}
\setbeamercolor{alertblock body}{bg=OtagoGold!8, fg=TextGrey}
\setbeamercolor{alerted text}{fg=OtagoGold!80!black}
\setbeamercolor{structure}{fg=OtagoBlue}

\setbeamertemplate{frametitle}{%
	\vspace{0.35em}%
	\textbf{\insertframetitle}\par%
	\vspace{0.05em}%
	\textcolor{OtagoGold}{\rule{\textwidth}{0.6pt}}}

\setbeamertemplate{navigation symbols}{}
\setbeamertemplate{footline}{%
	\leavevmode\hbox{%
		\begin{beamercolorbox}[wd=.40\paperwidth,ht=2.5ex,dp=1.5ex,leftskip=.3cm]{author in head/foot}%
			\textcolor{TextGrey}{\footnotesize ECON306 \textbullet\ University of Otago}%
		\end{beamercolorbox}%
		\begin{beamercolorbox}[wd=.47\paperwidth,ht=2.5ex,dp=1.5ex,center]{title in head/foot}%
			\textcolor{TextGrey}{\footnotesize \insertshorttitle}%
		\end{beamercolorbox}%
		\begin{beamercolorbox}[wd=.13\paperwidth,ht=2.5ex,dp=1.5ex,rightskip=.3cm,right]{date in head/foot}%
			\textcolor{TextGrey}{\footnotesize \insertframenumber\,/\,\inserttotalframenumber}%
	\end{beamercolorbox}}\vskip0pt}

\setbeamertemplate{itemize item}{\textcolor{OtagoGold}{\small$\blacktriangleright$}}
\setbeamertemplate{itemize subitem}{\textcolor{AccentBlue}{\small$\triangleright$}}
\setbeamertemplate{enumerate item}{\textcolor{OtagoBlue}{\small\bfseries\insertenumlabel.}}

\setbeamertemplate{title page}{%
	\begin{tikzpicture}[remember picture,overlay]
		\fill[OtagoBlue] (current page.north west)
		rectangle ([yshift=-0.38\paperheight]current page.north east);
	\end{tikzpicture}
	\vspace{0.06\paperheight}
	\begin{beamercolorbox}[wd=\textwidth, sep=0pt]{title}
		{\color{white}\Large\textbf{\inserttitle}}\\[0.4em]
		{\color{OtagoGold}\large\insertsubtitle}
	\end{beamercolorbox}
	\vspace{1.4em}
	\begin{beamercolorbox}[wd=\textwidth, sep=0pt]{author}
		{\color{TextGrey}\normalsize\insertauthor}\\[0.2em]
		{\color{TextGrey}\small\insertinstitute}\\[0.2em]
		{\color{TextGrey}\small\insertdate}
\end{beamercolorbox}}

% ─── Packages ──────────────────────────────────────────────
\usepackage{tikz}
\usepackage{booktabs}
\usepackage{array}
\usepackage{graphicx}
\usepackage{hyperref}
\usepackage{amsmath}
\usepackage{amssymb}
\usepackage{colortbl}
\usepackage{multirow}
\usepackage{xcolor}
\hypersetup{colorlinks=true, linkcolor=AccentBlue, urlcolor=AccentBlue,
	citecolor=AccentBlue, pdfauthor={Austin Henderson},
	pdftitle={ECON306 S2 2026}}

% ─── Images folder ─────────────────────────────────────────
%  Place all slide images in the 'images/' subfolder next to each deck.
%  Usage: \includegraphics[width=0.85\textwidth]{figure_name}
\graphicspath{{images/}}

% ─── Reusable macros ───────────────────────────────────────
% Recap slide at start of each lecture
\newcommand{\recapslide}[1]{%
	\begin{frame}{Recap: What Did We Cover Last Time?}
		\begin{block}{Key points from last lecture}
			\begin{itemize}#1\end{itemize}
		\end{block}
\end{frame}}

% Plan slide at start of each lecture
\newcommand{\planslide}[2]{%
	\begin{frame}{Today's Plan — #1}
		\begin{columns}
			\column{0.6\textwidth}
			\begin{itemize}#2\end{itemize}
			\column{0.38\textwidth}
			\begin{alertblock}{Reading}
				\footnotesize See Brightspace (Aoroa) for required readings before this lecture.
			\end{alertblock}
		\end{columns}
\end{frame}}

% Summary slide at end of each lecture
\newcommand{\summaryside}[1]{%
	\begin{frame}{Lecture Summary}
		\begin{exampleblock}{Key takeaways}
			\begin{itemize}#1\end{itemize}
		\end{exampleblock}
		\vspace{0.5em}
		{\footnotesize\textcolor{OtagoGold}{\textbf{Brightspace (Aoroa):}} slides, readings, and any announcements posted there.}
\end{frame}}

% NZ spotlight box
\newcommand{\nzbox}[2]{%
	\begin{alertblock}{\includegraphics[height=0.8em]{nz_flag_icon}\hspace{0.3em} NZ Spotlight: #1}
		#2
\end{alertblock}}

% Tutorial callout
\newcommand{\tutorialbox}[1]{%
	\begin{block}{\textbf{Tutorial This Week}}
		#1
\end{block}}

% ─── Section dividers ──────────────────────────────────────
\AtBeginSection[]{
	\begin{frame}[plain]
		\begin{tikzpicture}[remember picture,overlay]
			\fill[OtagoBlue!95] (current page.north west)
			rectangle (current page.south east);
			\node[anchor=center, text=white, font=\Large\bfseries,
			text width=0.9\paperwidth, align=center]
			at (current page.center) {\insertsectionhead};
			\node[anchor=south, text=OtagoGold, font=\normalsize,
			text width=0.9\paperwidth, align=center]
			at ([yshift=2.5cm]current page.south) {ECON306 — University of Otago};
		\end{tikzpicture}
\end{frame}}

%  Do NOT include \documentclass or per-deck metadata here.
% ============================================================

% ─── Theme & Colours ────────────────────────────────────────
\usetheme{default}
\usecolortheme{default}
\usefonttheme{professionalfonts}
\definecolor{OtagoBlue}{HTML}{003057}
\definecolor{OtagoGold}{HTML}{A8923A}
\definecolor{TextGrey}{HTML}{3A3A3A}
\definecolor{AccentBlue}{HTML}{2B6CB0}
\definecolor{LightBg}{HTML}{F7F9FC}

\setbeamercolor{frametitle}{fg=OtagoBlue, bg=white}
\setbeamercolor{title}{fg=OtagoBlue}
\setbeamercolor{subtitle}{fg=OtagoGold}
\setbeamercolor{author}{fg=TextGrey}
\setbeamercolor{date}{fg=TextGrey}
\setbeamercolor{institute}{fg=TextGrey}
\setbeamercolor{normal text}{fg=TextGrey, bg=white}
\setbeamercolor{item}{fg=OtagoBlue}
\setbeamercolor{subitem}{fg=AccentBlue}
\setbeamercolor{block title}{bg=OtagoBlue!12, fg=OtagoBlue}
\setbeamercolor{block body}{bg=OtagoBlue!5, fg=TextGrey}
\setbeamercolor{alertblock title}{bg=OtagoGold!20, fg=OtagoBlue}
\setbeamercolor{alertblock body}{bg=OtagoGold!8, fg=TextGrey}
\setbeamercolor{alerted text}{fg=OtagoGold!80!black}
\setbeamercolor{structure}{fg=OtagoBlue}

\setbeamertemplate{frametitle}{%
	\vspace{0.35em}%
	\textbf{\insertframetitle}\par%
	\vspace{0.05em}%
	\textcolor{OtagoGold}{\rule{\textwidth}{0.6pt}}}

\setbeamertemplate{navigation symbols}{}
\setbeamertemplate{footline}{%
	\leavevmode\hbox{%
		\begin{beamercolorbox}[wd=.40\paperwidth,ht=2.5ex,dp=1.5ex,leftskip=.3cm]{author in head/foot}%
			\textcolor{TextGrey}{\footnotesize ECON306 \textbullet\ University of Otago}%
		\end{beamercolorbox}%
		\begin{beamercolorbox}[wd=.47\paperwidth,ht=2.5ex,dp=1.5ex,center]{title in head/foot}%
			\textcolor{TextGrey}{\footnotesize \insertshorttitle}%
		\end{beamercolorbox}%
		\begin{beamercolorbox}[wd=.13\paperwidth,ht=2.5ex,dp=1.5ex,rightskip=.3cm,right]{date in head/foot}%
			\textcolor{TextGrey}{\footnotesize \insertframenumber\,/\,\inserttotalframenumber}%
	\end{beamercolorbox}}\vskip0pt}

\setbeamertemplate{itemize item}{\textcolor{OtagoGold}{\small$\blacktriangleright$}}
\setbeamertemplate{itemize subitem}{\textcolor{AccentBlue}{\small$\triangleright$}}
\setbeamertemplate{enumerate item}{\textcolor{OtagoBlue}{\small\bfseries\insertenumlabel.}}

\setbeamertemplate{title page}{%
	\begin{tikzpicture}[remember picture,overlay]
		\fill[OtagoBlue] (current page.north west)
		rectangle ([yshift=-0.38\paperheight]current page.north east);
	\end{tikzpicture}
	\vspace{0.06\paperheight}
	\begin{beamercolorbox}[wd=\textwidth, sep=0pt]{title}
		{\color{white}\Large\textbf{\inserttitle}}\\[0.4em]
		{\color{OtagoGold}\large\insertsubtitle}
	\end{beamercolorbox}
	\vspace{1.4em}
	\begin{beamercolorbox}[wd=\textwidth, sep=0pt]{author}
		{\color{TextGrey}\normalsize\insertauthor}\\[0.2em]
		{\color{TextGrey}\small\insertinstitute}\\[0.2em]
		{\color{TextGrey}\small\insertdate}
\end{beamercolorbox}}

% ─── Packages ──────────────────────────────────────────────
\usepackage{tikz}
\usepackage{booktabs}
\usepackage{array}
\usepackage{graphicx}
\usepackage{hyperref}
\usepackage{amsmath}
\usepackage{amssymb}
\usepackage{colortbl}
\usepackage{multirow}
\usepackage{xcolor}
\hypersetup{colorlinks=true, linkcolor=AccentBlue, urlcolor=AccentBlue,
	citecolor=AccentBlue, pdfauthor={Austin Henderson},
	pdftitle={ECON306 S2 2026}}

% ─── Images folder ─────────────────────────────────────────
%  Place all slide images in the 'images/' subfolder next to each deck.
%  Usage: \includegraphics[width=0.85\textwidth]{figure_name}
\graphicspath{{images/}}

% ─── Reusable macros ───────────────────────────────────────
% Recap slide at start of each lecture
\newcommand{\recapslide}[1]{%
	\begin{frame}{Recap: What Did We Cover Last Time?}
		\begin{block}{Key points from last lecture}
			\begin{itemize}#1\end{itemize}
		\end{block}
\end{frame}}

% Plan slide at start of each lecture
\newcommand{\planslide}[2]{%
	\begin{frame}{Today's Plan — #1}
		\begin{columns}
			\column{0.6\textwidth}
			\begin{itemize}#2\end{itemize}
			\column{0.38\textwidth}
			\begin{alertblock}{Reading}
				\footnotesize See Brightspace (Aoroa) for required readings before this lecture.
			\end{alertblock}
		\end{columns}
\end{frame}}

% Summary slide at end of each lecture
\newcommand{\summaryside}[1]{%
	\begin{frame}{Lecture Summary}
		\begin{exampleblock}{Key takeaways}
			\begin{itemize}#1\end{itemize}
		\end{exampleblock}
		\vspace{0.5em}
		{\footnotesize\textcolor{OtagoGold}{\textbf{Brightspace (Aoroa):}} slides, readings, and any announcements posted there.}
\end{frame}}

% NZ spotlight box
\newcommand{\nzbox}[2]{%
	\begin{alertblock}{\includegraphics[height=0.8em]{nz_flag_icon}\hspace{0.3em} NZ Spotlight: #1}
		#2
\end{alertblock}}

% Tutorial callout
\newcommand{\tutorialbox}[1]{%
	\begin{block}{\textbf{Tutorial This Week}}
		#1
\end{block}}

% ─── Section dividers ──────────────────────────────────────
\AtBeginSection[]{
	\begin{frame}[plain]
		\begin{tikzpicture}[remember picture,overlay]
			\fill[OtagoBlue!95] (current page.north west)
			rectangle (current page.south east);
			\node[anchor=center, text=white, font=\Large\bfseries,
			text width=0.9\paperwidth, align=center]
			at (current page.center) {\insertsectionhead};
			\node[anchor=south, text=OtagoGold, font=\normalsize,
			text width=0.9\paperwidth, align=center]
			at ([yshift=2.5cm]current page.south) {ECON306 — University of Otago};
		\end{tikzpicture}
\end{frame}}


\title[ECON306: W11D2 — HC vs. Signalling Evidence]{ECON306\\The Economics of Health and Education}
\subtitle{W11D2: Human Capital vs.\ Signalling: Evidence; Sheepskin Effects}
\author{Austin Henderson}
\institute{Department of Economics\\University of Otago}
\date{Thursday 1 October 2026}

\begin{document}

\begin{frame}[plain]
  \titlepage
\end{frame}

\recapslide{
  \item Spence (1973): education may signal ability without raising productivity --- social returns potentially zero
  \item Separating equilibrium: high-ability workers signal; low-ability do not (cost differential is the key)
  \item Single-crossing property must hold for separating equilibrium to be stable
}

\planslide{Human Capital vs.\ Signalling: What Does the Evidence Say?}{
  \item Empirical strategies for distinguishing HC from signalling
  \item Evidence for human capital theory: natural experiments
  \item Evidence for signalling: sheepskin effects and self-employed wage returns
  \item The consensus: probably both matter --- but how much of each?
  \item Policy implications of the balance of evidence
}

\begin{frame}{Why Is This Hard to Test?}
  \begin{columns}
    \column{0.52\textwidth}
      \small \textbf{The identification problem:} Both theories predict the same stylised fact
      --- graduates earn more. To distinguish them, we need variation in education
      \textbf{unrelated to ability}.
      \begin{itemize}
        \item \textbf{Ability bias}: more able people earn more \emph{and} are more likely
              to attend university
        \item \textbf{Selection bias}: choosers differ systematically from non-choosers
        \item Both biases \emph{overstate} returns in naive OLS
      \end{itemize}
    \column{0.45\textwidth}
      \begin{alertblock}{Solution strategies}
        \small
        \begin{itemize}
          \item \textbf{IV/natural experiments}: compulsory schooling law changes (Angrist \& Krueger, 1991)
          \item \textbf{Twin studies}: identical twins with different education levels
          \item \textbf{Sheepskin tests}: wages at degree completion milestones
          \item \textbf{Self-employed comparison}: no employer to signal to
        \end{itemize}
      \end{alertblock}
  \end{columns}
\end{frame}

\begin{frame}{Evidence for Human Capital Theory}
  \begin{columns}
    \column{0.52\textwidth}
      \footnotesize \textbf{Natural experiments (IV):}
      \begin{itemize}
        \item \textbf{Angrist \& Krueger (1991)}: quarter of birth as instrument for
              years of schooling; extra schooling raises wages --- suggests real productivity gain
        \item Similar results from compulsory schooling law changes in UK, Denmark, Sweden, NZ
      \end{itemize}
      \textbf{Vocational training:}
      \begin{itemize}
        \item Apprenticeships and vocational training produce wage gains even where there is no
              signalling component
      \end{itemize}
    \column{0.45\textwidth}
      \begin{exampleblock}{NZ evidence}
        \footnotesize Using the NZ Longitudinal Business Database (LBD), extra years of schooling
        raise wages even after controlling for ability proxies.\\[0.2em]
        Returns persist for occupations where degrees are not formally required for entry ---
        hard to explain under pure signalling.
      \end{exampleblock}
  \end{columns}
\end{frame}

\begin{frame}{Evidence for Signalling: Sheepskin Effects}
  \begin{columns}
    \column{0.55\textwidth}
      \begin{block}{The sheepskin effect}
        \footnotesize ``Sheepskin'' = degree certificate (historically on sheepskin).\\[0.1em]
        \textbf{Key finding}: wages jump \textbf{discontinuously} at degree completion,
        beyond what an extra year of study predicts under HC theory.
        \begin{itemize}
          \item Years 1--3: small incremental wage gains
          \item Year 4 (graduation): disproportionately large wage jump
        \end{itemize}
      \end{block}
      \footnotesize More consistent with \textbf{signalling} than pure HC (skills accumulate gradually, credentials do not).
    \column{0.42\textwidth}
      \begin{alertblock}{NZ and international evidence}
        \footnotesize Sheepskin effects documented in NZ, US, UK, and Australia. Effect size
        typically 2--4 \textit{extra} years of education. Consistent with credential-based
        signalling.
      \end{alertblock}
  \end{columns}
\end{frame}

\begin{frame}{Evidence for Signalling: Self-Employed Workers}
  \begin{columns}
    \column{0.52\textwidth}
      \begin{alertblock}{Key test}
        \small Under signalling theory:
        \begin{itemize}
          \item Education signals ability to \textbf{employers}
          \item Self-employed have no employer to signal to
          \item $\Rightarrow$ Returns should be \textbf{lower for the self-employed}
        \end{itemize}
      \end{alertblock}
    \column{0.45\textwidth}
      \begin{block}{Evidence}
        \small US, UK, and NZ: returns to schooling are \textbf{lower for self-employed} than
        employees.\\[0.2em]
        Difference is \textbf{partial}, not complete --- consistent with a mixed HC +
        signalling story.
      \end{block}
  \end{columns}
  \vspace{0.1em}
  \begin{exampleblock}{Bryan Caplan (2018): \textit{The Case Against Education}}
    \footnotesize Argues signalling accounts for $\approx$80\% of the wage premium ---
    provocative but an important corrective to naive HC theory.
  \end{exampleblock}
\end{frame}

\begin{frame}{The Balance of Evidence}
  \begin{columns}
    \column{0.48\textwidth}
      \begin{block}{Evidence favouring human capital}
        \footnotesize
        \begin{itemize}
          \item IV estimates (Angrist-Krueger, compulsory schooling laws)
          \item Vocational training wage gains
          \item Returns persist in occupations with no formal credential requirement
          \item Employer-reported skill shortages in graduate fields
        \end{itemize}
      \end{block}
    \column{0.48\textwidth}
      \begin{alertblock}{Evidence favouring signalling}
        \footnotesize
        \begin{itemize}
          \item Sheepskin effects (wage jump at graduation)
          \item Lower returns for self-employed
          \item ``Signalling inflation'' in credential requirements over time
          \item Credential screening even for non-degree tasks
        \end{itemize}
      \end{alertblock}
  \end{columns}
  \vspace{0.3em}
  \footnotesize \textbf{The likely answer: both matter.} Importance varies by level (postgraduate $\Rightarrow$ more HC) and field (medicine/engineering $\Rightarrow$ more HC; arts $\Rightarrow$ more signalling).
\end{frame}

\begin{frame}{Implications for NZ Education Policy}
  \small
  \begin{enumerate}
    \item \textbf{If HC dominates}: strong case for public subsidy (positive externalities +
          positive social returns $\Rightarrow$ under-investment without subsidy)
    \item \textbf{If signalling dominates}: subsidising education accelerates the signalling
          arms race; reducing subsidies and moving towards user-pays may be more efficient
    \item \textbf{If both matter (likely)}: some subsidy is justified to internalise
          externalities; income-contingent loans solve the capital market failure component;
          exact mix depends on empirical estimates of social returns
    \item \textbf{The level matters}: primary and secondary education is more human capital;
          graduate/postgraduate has a larger signalling component
          $\Rightarrow$ tapering subsidy rates with level may be efficient
  \end{enumerate}
\end{frame}

\summaryside{
  \footnotesize
  \item Identifying the HC--signalling split requires exogenous variation in education (IV, natural experiments, twin studies, self-employment)
  \item Evidence for HC: compulsory schooling IV (Angrist \& Krueger, 1991); vocational training wage gains; NZ LBD results
  \item Evidence for signalling: sheepskin effects (discontinuous wage jump at graduation); lower returns for self-employed
  \item Likely answer: both matter; relative importance varies by education level and field; policy subsidy should reflect net social returns
}

\end{document}
